\chapter{Why Supersymmetry Might Matter}

\lectureoverview{Vacuum fluctuations provide the thread that joins the apparently separate subjects of this chapter. Boundary conditions turn vacuum fluctuations into the Casimir force; loop effects threaten the Higgs mass; new virtual particles alter the running of gauge couplings; and vacuum energy eventually controls the expansion of the universe. These problems lead to three practical motivations for low-energy supersymmetry: stabilizing the Higgs scale, supplying a dark-matter candidate, and improving gauge-coupling unification. We then return to the formal construction. Supersymmetric multiplets become superfields on a space enlarged by Grassmann coordinates, and Berezin integration extracts the top component needed for a supersymmetric action.}

\section{Vacuum Fluctuations Between Reflecting Plates}

We return to supersymmetry by first settling a question about the Casimir effect. Place two parallel reflecting plates a distance \(d\) apart. Even when no real particles are present, a quantum field has vacuum fluctuations. We may picture them as virtual processes, or equivalently as the zero-point motion of the field's oscillator modes. The plates impose boundary conditions, so the allowed fluctuations between them differ from those in unrestricted space.

The important quantity is not an absolute zero of energy but the change produced when the separation changes. Denote the vacuum energy per unit area by \(\mathcal E(d)=E(d)/A\). Once \(\mathcal E\) depends on \(d\), it is a potential energy, and the force per unit area is
\begin{equation}
\frac{F}{A}=-\frac{d\mathcal E}{dd}.
\end{equation}
Thus a rearrangement of virtual processes becomes an observable force. In the familiar electromagnetic example with ideal conducting plates, the energy and pressure are
\begin{align}
\mathcal E(d)&=-\frac{\pi^2}{720\,d^3},\\
\frac{F}{A}&=-\frac{\pi^2}{240\,d^4},
\end{align}
in units \(\hbar=c=1\).\editorialnote{These coefficients and the distinction between the \(d^{-3}\) energy per area and the \(d^{-4}\) pressure resolve the dimensional question raised at the board. They are the standard ideal-plate result, not a derivation carried out here.} The minus sign denotes attraction.

\begin{classroomqa}
\audiencequestion{Does the fourth power describe the force or the energy? \lecturetimestamp{00:04:39}}
\lecturerresponse{For parallel plates in four spacetime dimensions, the energy per unit area falls as \(d^{-3}\), while differentiating with respect to \(d\) makes the force per unit area fall as \(d^{-4}\). Dimensional analysis already fixes these powers once \(\hbar\) and \(c\) are restored.}
\end{classroomqa}

One sometimes summarizes the calculation by saying that the plates interrupt virtual photon loops. That picture is useful provided it is not taken too literally. The precise calculation is a change in the spectrum of allowed field modes or, equivalently, a change in the quantum effective action. Either language leads to the same separation-dependent energy.

\begin{classroomqa}
\audiencequestion{Can the Casimir force be repulsive rather than attractive? \lecturetimestamp{00:07:48}}
\lecturerresponse{It can. The sign is not universal; it depends on the fields, geometry, materials, and boundary conditions. Fermionic loops often bring the opposite loop sign from bosonic loops, which motivates the suggestion made at the board, but fermionic statistics by itself is not a general theorem that every Casimir force must be repulsive.\editorialnote{This qualification replaces a deliberately tentative classroom answer with the general statement needed to avoid a false universal rule.}}
\end{classroomqa}

\subsection{What would count as evidence for vacuum energy?}

The Casimir force is measured, and atomic observables are sensitive to vacuum loops as well. Yet neither fact by itself demonstrates that an absolute vacuum energy gravitates. A Casimir experiment measures an energy difference caused by boundaries; the same force can be described without assigning operational meaning to an unbounded absolute zero-point sum.

The sharp test comes from gravity. Energy is gravitational charge. If empty space carries a uniform energy density, it must enter Einstein's equations and influence cosmic expansion. The accelerated expansion of the universe is therefore the direct gravitational evidence for a vacuum-like component.

\begin{classroomqa}
\audiencequestion{Why do popular accounts call vacuum energy mysterious, as if physicists do not know what it is? \lecturetimestamp{00:08:19}}
\lecturerresponse{Quantum field theory gives every field a zero-point contribution, so the possibility of vacuum energy is not mysterious. The mystery is quantitative: why is the gravitating value so extraordinarily small compared with estimates based on known quantum fields and a high-energy cutoff? The surprising fact is its near absence, not its bare presence.}
\end{classroomqa}

For many years the observational upper bound was so small that it was psychologically tempting to set the cosmological constant exactly to zero and search for a principle that enforced the cancellation. Improved cosmological observations instead found a small nonzero value. That discovery did not make vacuum energy conceptually alien; it made the unexplained degree of suppression impossible to ignore. The relevant comparison depends on the assumed ultraviolet scale, but every natural high-energy estimate demands an enormous cancellation.

\section{What One Supercharge Does to a Many-Particle State}

Before listing reasons to believe in supersymmetry, we should remove a possible paradox. Bosons may occupy one mode in arbitrarily large numbers, while a fermionic mode has occupation zero or one. If a symmetry exchanges bosons and fermions, does it turn a thousand bosons in one mode into a thousand forbidden fermions?

No. A fermionic generator changes one quantum at a time. In the oscillator model, let \(a\) annihilate a boson and \(c^\dagger\) create its fermionic partner. One supercharge may be represented as
\begin{equation}
Q=c^\dagger a.
\end{equation}
On a state with \(n\) bosons and no fermion,
\begin{equation}
Q\ket{n,0}
=\sqrt n\,\ket{n-1,1}.
\end{equation}
One boson has been removed and one fermion inserted. Acting again with the same charge gives zero:
\begin{equation}
Q^2=(c^\dagger)^2a^2=0.
\end{equation}
This is the nilpotence inherited from fermionic creation operators.

\begin{classroomqa}
\audiencequestion{If many bosons occupy the same state, why does a supersymmetry transformation not produce many fermions in that same state and violate the exclusion principle? \lecturetimestamp{00:12:05}}
\lecturerresponse{The generator does not replace the entire population. It changes the state by one unit of fermion number. The same fermionic charge cannot be compounded like an ordinary rotation generator because its square vanishes. A complementary charge can undo the exchange, and the mixed anticommutator closes on time translation.}
\end{classroomqa}

The last sentence is essential. The shorthand \(Q^2=0\) applies to an individual chiral component or to the simple charge just written; it does not mean that every product of supercharges vanishes. In four-dimensional supersymmetry,
\begin{equation}
\{Q_\alpha,\bar Q_{\dot\beta}\}
=2\sigma^\mu_{\alpha\dot\beta}P_\mu,
\end{equation}
while
\begin{equation}
\{Q_\alpha,Q_\beta\}=0,
\qquad
\{\bar Q_{\dot\alpha},\bar Q_{\dot\beta}\}=0
\end{equation}
for the minimal algebra without central charges.\editorialnote{This component form makes precise the classroom shorthand that repeated action of the same \(Q\) vanishes while a supercharge and its conjugate generate translation.}

\section{Three Reasons to Look for Low-Energy Supersymmetry}

\begin{classroomqa}
\audiencequestion{Is there motivation for supersymmetry besides cancelling vacuum energy? \lecturetimestamp{00:15:06}}
\lecturerresponse{There are three closely related motivations: control of the Higgs self-energy, a possible stable dark-matter particle, and a much sharper meeting of the gauge couplings at high energy. None proves that supersymmetry is realized in nature, but together they make it a serious theory to test.}
\end{classroomqa}

The phrase \emph{low-energy supersymmetry} is relative. It does not mean that superpartners are light on an everyday scale; it means that their masses are not vastly above the electroweak and accelerator scales. All three arguments weaken if the superpartners are pushed arbitrarily high.

\subsection{Protecting the Higgs scale}

A scalar mass is unusually sensitive to high-energy physics. Schematically, a bosonic loop produces a correction of the form
\begin{equation}
\delta m_H^2\big|_{\mathrm B}
\sim +\frac{\lambda}{16\pi^2}\Lambda^2,
\end{equation}
where \(\Lambda\) denotes the scale up to which the effective theory is being used. A fermionic loop carries the opposite sign,
\begin{equation}
\delta m_H^2\big|_{\mathrm F}
\sim -\frac{|y|^2}{8\pi^2}\Lambda^2.
\end{equation}
The coefficients depend on the interaction and field multiplicities, but supersymmetry relates them so that the leading ultraviolet sensitivity cancels between partners.\editorialnote{The schematic loop formulas expose the mathematics behind the qualitative self-energy argument. They are not intended as a complete Standard Model Higgs-mass calculation.}

This matters because the Higgs expectation value sets the electroweak scale. Gauge-boson and fermion masses obey relations such as
\begin{equation}
m_W=\frac{gv}{2},
\qquad
m_f=\frac{y_fv}{\sqrt2}.
\end{equation}
If the Higgs sector were naturally driven to an enormous scale, the familiar low-energy particle spectrum would be lost unless bare parameters were adjusted against radiative corrections with extreme precision.

Exact supersymmetry would make the cancellation exact, but exact supersymmetry is already excluded by the unequal masses of observed bosons and fermions. With softly broken supersymmetry, the quadratic sensitivity is replaced by dependence on superpartner splittings, schematically
\begin{equation}
\delta m_H^2
\sim
\frac{g^2}{16\pi^2}
\left(m_{\mathrm{soft}}^2-m_{\mathrm{ordinary}}^2\right)
\ln\frac{\Lambda^2}{m_{\mathrm{soft}}^2}.
\end{equation}
The cancellation remains useful only while the splitting is not too large.

\subsection{A stable endpoint for dark-matter cascades}

Separate the spectrum into ordinary particles and their superpartners. A fermion has a bosonic partner; a boson has a fermionic partner. If the interactions preserve a parity that distinguishes these two classes, a superpartner decay must leave another superpartner:
\begin{equation}
\widetilde X
\longrightarrow X+\widetilde Y
\longrightarrow X+Y+\widetilde Z
\longrightarrow\cdots.
\end{equation}
The cascade cannot discard its last superpartner. It therefore ends at the lightest supersymmetric particle, the LSP.

The commonly used symmetry is \(R\)-parity,
\begin{equation}
R_p=(-1)^{3(B-L)+2s}.
\end{equation}
Standard Model particles have \(R_p=+1\), while their superpartners have \(R_p=-1\). If \(R_p\) is exact, every interaction contains an even number of superpartners and the LSP is stable.\editorialnote{Supersymmetry alone does not guarantee a stable LSP. Stability requires conserved \(R\)-parity or another symmetry with the same consequence; without it, the LSP may decay.}

Dark matter must survive for at least cosmological times. It must also interact weakly enough to pass repeatedly through the visible galactic disk without losing substantial energy. An electrically charged or strongly interacting population would scatter, dissipate energy, and tend to collapse toward the same thin structure as ordinary matter. The observed extended halo instead points toward a neutral, weakly interacting component.

This does not identify the LSP uniquely. Depending on the model and symmetry-breaking pattern, the candidate might be a neutral mixture of electroweak gauginos and Higgsinos, or it might be a gravitino. Charged or colored candidates are strongly constrained. What supersymmetry supplies is not a guaranteed answer but a natural mechanism for producing a neutral, long-lived endpoint.

\section{Running Couplings and the Hint of Unification}

The third motivation starts from charge screening. A measured coupling depends on the energy scale because virtual particles dress the interaction. For electromagnetism, an electron-positron loop inserted in a photon propagator screens the charge at long distance. At shorter distance, or higher momentum transfer, less of that polarization cloud lies between the probes and the effective charge is larger.

Write
\begin{equation}
\alpha_i=\frac{g_i^2}{4\pi}.
\end{equation}
At one loop the inverse coupling runs linearly with the logarithm of energy:
\begin{equation}
\frac{d\alpha_i^{-1}}{d\ln\mu}
=-\frac{b_i}{2\pi}.
\end{equation}
Thus
\begin{equation}
\alpha_i^{-1}(\mu)
=\alpha_i^{-1}(\mu_0)
-\frac{b_i}{2\pi}\ln\frac{\mu}{\mu_0}.
\end{equation}
\editorialnote{This one-loop equation is the precise version of the nearly straight lines drawn as functions of logarithmic energy. Thresholds and higher-loop corrections bend the lines slightly.}

Above the electroweak scale the three fundamental gauge couplings are those of
\begin{equation}
U(1)_Y,\qquad SU(2)_L,\qquad SU(3)_c.
\end{equation}
The electromagnetic coupling used for intuition at low energy is a combination of the first two after electroweak symmetry breaking. For QCD the non-Abelian gauge fields produce antiscreening: the strong coupling decreases at short distance, the phenomenon of asymptotic freedom. Consequently \(\alpha_3^{-1}\) rises with \(\ln\mu\), while the other inverse couplings have different slopes.

The Standard Model curves point suggestively toward one another but do not meet at a single point. Superpartners modify the virtual loops and therefore change the coefficients \(b_i\). In the conventional grand-unified normalization,
\begin{align}
(b_1,b_2,b_3)_{\mathrm{SM}}
&=\left(\frac{41}{10},-\frac{19}{6},-7\right),\\
(b_1,b_2,b_3)_{\mathrm{MSSM}}
&=\left(\frac{33}{5},1,-3\right).
\end{align}
With superpartner thresholds near the TeV scale, the extrapolated supersymmetric curves meet much more accurately near
\begin{equation}
M_{\mathrm{GUT}}\sim 10^{16}\ \mathrm{GeV}.
\end{equation}
No experiment directly reaches that scale. Low-energy couplings are measured, the field content determines the running, and the curves are extrapolated over many decades. Two lines always cross somewhere; three independent lines meeting near one point is the nontrivial observation.

\begin{classroomqa}
\audiencequestion{Does the improved extrapolation assume exact supersymmetry and equal partner masses? \lecturetimestamp{00:38:09}}
\lecturerresponse{No. Supersymmetry is broken, so the masses are split. What matters is when the superpartner loops enter the running. If the partners are not too heavy, the slopes change early enough for the three couplings to converge. If they are exceedingly heavy, the supersymmetric thresholds arrive too late and the improved meeting is lost.}
\end{classroomqa}

Threshold corrections and the unknown high-energy spectrum prevent a literal claim that the three lines meet with mathematical exactness. The result is nevertheless striking: the same approximate superpartner scale that helps the Higgs and provides a thermal dark-matter candidate also improves unification.

\subsection{Gravity and the common high-energy scale}

Gravity behaves differently because its charge is energy itself. A useful dimensionless measure at center-of-mass energy \(E\) is
\begin{equation}
\alpha_G(E)\sim GE^2
=\left(\frac{E}{M_{\mathrm{Pl}}}\right)^2.
\end{equation}
At ordinary energies this is fantastically small, but it grows quadratically with \(E\). Its inverse therefore falls much faster than the logarithmically running gauge inverse couplings and becomes order unity near the Planck scale.

This comparison suggests that all interactions may be governed by neighboring high scales, but it is less clean than gauge unification. The gauge couplings are dimensionless parameters in renormalizable field theories; Newton's constant is dimensionful, so translating gravity onto the same plot requires a convention.

\begin{classroomqa}
\audiencequestion{Do superpartners participate in the same three gauge interactions, or do they define new forces? \lecturetimestamp{00:42:35}}
\lecturerresponse{They carry the same gauge quantum numbers as their partners. A selectron has the electron's electric charge, and a squark has the corresponding quark's color and electroweak charges, though their spins and masses differ. Their loops therefore modify the same gauge interactions rather than introducing three new ones.}
\end{classroomqa}

\begin{classroomqa}
\audiencequestion{Would discovery of a superpartner at the LHC be definitive evidence for supersymmetry? \lecturetimestamp{00:43:27}}
\lecturerresponse{A single unfamiliar particle would not be enough by its name alone, but a superpartner spectrum obeying supersymmetric relations among spins, charges, masses, and couplings would be highly distinctive. Supersymmetry imposes many correlated restrictions, so the detailed interaction pattern would provide the decisive test.}
\end{classroomqa}

\subsection{Why the scale cannot be moved away without cost}

If superpartners lie only a little beyond experimental reach, all three motivations can still operate. If their masses are raised by orders of magnitude, each argument deteriorates:
\begin{enumerate}
\item Higgs-mass cancellations leave a large residual sensitivity to the splitting.
\item A conventional weak-scale thermal relic no longer emerges at the expected scale.
\item Gauge-running thresholds move upward and the three couplings meet less accurately.
\end{enumerate}
Supersymmetry could still exist at a much higher scale, but it would no longer solve these particular low-energy problems in the same economical way. Longevity of the idea is not evidence that nature must use it; it is evidence that the proposal is coherent enough to deserve a decisive test.

\section{What Cosmology Can and Cannot Tell Us About Dark Matter}

\begin{classroomqa}
\audiencequestion{Can standard cosmology determine the amount of dark matter and the mass of each dark-matter particle, much as nucleosynthesis constrains hydrogen and helium? \lecturetimestamp{00:46:32}}
\lecturerresponse{Cosmology measures the total dark-matter mass density very well, but density alone does not separate mass from number. The same density could be carried by a few heavy particles or an enormous number of very light ones. Additional theory or nongravitational observations are needed to identify the microscopic particle.}
\end{classroomqa}

Two broad possibilities illustrate the ambiguity. A weakly interacting particle near the electroweak or TeV scale can freeze out from a hot thermal bath. At the other extreme, an extremely light boson such as an axion can occupy a coherent field configuration or condensate. Intermediate possibilities are not forbidden merely because the two best-developed mechanisms sit near opposite ends.

The halo structure itself supplies an important constraint. Dark matter crosses ordinary matter with little friction, so it remains spatially extended instead of collapsing into the luminous disk. Rotation curves and other gravitational phenomena tell us that the mass is present, but they do not by themselves count the particles.

\begin{classroomqa}
\audiencequestion{Could the early universe initially have contained only superparticles, which later decayed into the particles we see? \lecturetimestamp{00:49:11}}
\lecturerresponse{At temperatures much larger than all relevant masses, rapid reactions drive the species toward thermal equilibrium rather than an arbitrary one-sided population. Ordinary collisions create superpartners when they are underabundant, while superpartner collisions and decays replenish ordinary particles. The precise abundance depends on the interactions and degrees of freedom, but the hot plasma continually tends toward a balance.}
\end{classroomqa}

As the universe cools, massive species become Boltzmann suppressed and annihilate or decay. A stable LSP ceases to track equilibrium when its interaction rate becomes slower than cosmic expansion:
\begin{equation}
\Gamma_{\mathrm{ann}}
\sim n\langle\sigma v\rangle
\lesssim H.
\end{equation}
The remaining particles form a relic abundance.\editorialnote{This freeze-out criterion is the standard quantitative completion of the thermal-balance argument. Other LSP candidates, especially very weakly coupled gravitinos, can have nonthermal histories.}

\begin{classroomqa}
\audiencequestion{Which ordinary particle is the LSP the partner of, and must it be noninteracting? \lecturetimestamp{00:51:05}}
\lecturerresponse{There is no model-independent answer. A viable dark-matter LSP is expected to be electrically neutral and colorless, but it need not be completely noninteracting. It may be a mixture related to the photon, weak gauge bosons, and Higgs fields, or it may be the gravitino. Its precise identity follows from the supersymmetry-breaking spectrum.}
\end{classroomqa}

\section{Vacuum Diagrams, Self-Energy Diagrams, and Expansion}

\begin{classroomqa}
\audiencequestion{What is the difference between the Higgs divergence and the vacuum-energy divergence? \lecturetimestamp{00:52:09}}
\lecturerresponse{A vacuum diagram has no external particle lines; it changes the effective action of empty space and contributes to the vacuum energy. A Higgs self-energy diagram has two external Higgs legs; it changes the Higgs propagator and therefore its mass. Both involve virtual fluctuations, but they answer different questions.}
\end{classroomqa}

One may picture a real particle as an impurity inserted into the vacuum. Vacuum fluctuations that attach to this impurity shift its propagation and mass. The Casimir plates also modify vacuum fluctuations, but through macroscopic boundary conditions. The common ingredient is a response of quantum fluctuations to what has been inserted; the observable being corrected differs in each case.

\vspace{-2pt}
\subsection{Matter, radiation, and vacuum dilute differently}

Let \(a(t)\) be the cosmological scale factor and
\begin{equation}
H=\frac{\dot a}{a}
\end{equation}
the Hubble parameter. Local stress-energy conservation in a homogeneous universe gives
\begin{equation}
\dot\rho+3H(\rho+p)=0.
\end{equation}
For an equation of state \(p=w\rho\), this integrates to
\begin{equation}
\rho(a)\propto a^{-3(1+w)}.
\end{equation}
\editorialnote{The continuity equation supplies the compact derivation underlying the particle-counting and wavelength-stretching explanations.}

For cold, nonrelativistic matter, \(p\approx0\), so
\begin{equation}
\rho_m\propto a^{-3}.
\end{equation}
A fixed comoving volume grows as \(a^3\), while the rest energy carried by its particles remains approximately fixed.

For radiation, \(p=\rho/3\), and
\begin{equation}
\rho_r\propto a^{-4}.
\end{equation}
Three powers come from dilution of the number density. The fourth comes from redshift:
\begin{equation}
\lambda\propto a,
\qquad
E_\gamma=\frac{2\pi\hbar c}{\lambda}\propto a^{-1}.
\end{equation}

\begin{classroomqa}
\audiencequestion{If a photon has a wavefunction, why may we speak of it as a particle rather than a wave when deriving the dilution law? \lecturetimestamp{00:58:21}}
\lecturerresponse{The distinction does not affect the scaling. In the particle description the number in a comoving volume is fixed and each quantum redshifts as \(a^{-1}\). In the classical-wave description the wavelength stretches and the field energy falls by the same amount. We can use whichever language makes the expansion easiest to see.}
\end{classroomqa}

Vacuum energy has \(p_\Lambda=-\rho_\Lambda\), so
\begin{equation}
\rho_\Lambda=\text{constant}.
\end{equation}
It does not dilute as space expands. Indeed, a comoving region contains
\begin{equation}
E_\Lambda=\rho_\Lambda a^3V_0,
\end{equation}
which grows with its physical volume. That fact creates the energy-conservation question we shall address shortly.

The different powers organize cosmic history. Radiation falls fastest, matter falls more slowly, and vacuum energy does not fall at all. A tiny vacuum term can therefore be negligible in the early universe and dominant at late times. Using the rough values discussed here, the present energy budget is about \(70\%\) vacuum energy and \(30\%\) matter, with radiation much smaller.

\section{Geometry Equals Energy}

The time-time component of Einstein's equations turns the preceding scaling laws into dynamics. The board begins with a slogan and a single geometric term.

\begin{figure}[H]
\centering
\includegraphics[width=0.68\textwidth]{lecture_06_figure_03.png}
\lectureframe{01:06:37.088}
\caption{The organizing slogan \emph{geometry equals energy}, with the expansion term \(H^2=(\dot a/a)^2\) written beneath it.}
\end{figure}

For a homogeneous and isotropic universe, the first Friedmann equation is
\begin{equation}
H^2+\frac{k}{a^2}
=\frac{8\pi G}{3}\rho,
\end{equation}
in units \(c=1\), where \(k=+1,0,-1\) denotes closed, flat, or open spatial geometry and
\begin{equation}
\rho
=\rho_{\mathrm{ordinary}}
+\rho_{\mathrm{dark}}
+\rho_{\mathrm{radiation}}
+\rho_{\mathrm{vacuum}}.
\end{equation}
The sign seen beside \(1/a^2\) depends on whether the curvature term is kept on the left or moved to the right, as well as on the value of \(k\).

\begin{figure}[H]
\centering
\includegraphics[width=0.96\textwidth]{lecture_06_figure_04.png}
\lectureframe{01:08:46.794}
\caption{The full board balance: expansion and spatial curvature on the geometric side, with matter, radiation, and vacuum energy on the source side. Matter is further separated into ordinary and dark components.}
\end{figure}

\subsection{The cosmological Hamiltonian constraint}

\begin{classroomqa}
\audiencequestion{If the vacuum energy in a comoving region grows as the universe expands, is there a constant net energy density or a conserved total energy? \lecturetimestamp{01:01:51}}
\lecturerresponse{General relativity does not supply a universal global energy for an arbitrary expanding spacetime. In the homogeneous cosmological reduction, however, the time-time Einstein equation is a Hamiltonian constraint. When geometry and matter are placed on one side, the constrained Hamiltonian vanishes. The changing matter terms are balanced by the gravitational degrees of freedom rather than by an ordinary fixed energy in a box.}
\end{classroomqa}

Schematically, the canonical statement is
\begin{equation}
\mathcal H_{\mathrm{gravity}}
+\mathcal H_{\mathrm{matter}}
=0.
\end{equation}
This is the controlled meaning of saying that the total energy is zero. It is an equation of motion enforced by time-reparameterization invariance, not a claim that every local energy density vanishes. In a comoving volume, ordinary matter energy is roughly fixed, radiation energy decreases, and vacuum energy grows; meanwhile the gravitational contribution and expansion rate adjust so that Einstein's constraint remains satisfied.

Local conservation remains exact:
\begin{equation}
\nabla_\mu T^{\mu\nu}=0.
\end{equation}
For the homogeneous fluid this is precisely the continuity equation
\(\dot\rho+3H(\rho+p)=0\). What generally fails is the flat-spacetime intuition that there must also be one coordinate-independent global sum of matter energy conserved in an expanding universe.\editorialnote{This distinction makes the classroom statement that energy is always zero precise: local covariant conservation and the cosmological Hamiltonian constraint are exact, while a unique global matter-plus-gravity energy is not generally defined.}

\begin{classroomqa}
\audiencequestion{Is the zero constraint the condition that decides whether the universe is open, closed, or flat? \lecturetimestamp{01:04:42}}
\lecturerresponse{No. The Hamiltonian constraint applies in every case. Open, flat, and closed universes differ through the curvature label \(k\) and therefore through the \(k/a^2\) term. Curvature changes one contribution to the balance; it does not turn the constraint on or off.}
\end{classroomqa}

\section{Grassmann Coordinates and Superfields}

\begin{classroomqa}
\audiencequestion{Do Grassmann numbers actually appear in the Lagrangian of the theory? \lecturetimestamp{01:09:37}}
\lecturerresponse{They appear first because fermion fields are Grassmann odd. They appear a second time as coordinates of superspace. Superfields depend on those odd coordinates, and supersymmetric actions are built by differentiating and integrating with respect to them.}
\end{classroomqa}

An ordinary field is a function of spacetime,
\begin{equation}
\phi=\phi(x^\mu).
\end{equation}
Minimal supersymmetry in four dimensions augments \(x^\mu\) by a two-component odd spinor \(\theta^\alpha\) and its conjugate \(\bar\theta^{\dot\alpha}\):
\begin{equation}
z^M=(x^\mu,\theta^\alpha,\bar\theta^{\dot\alpha}).
\end{equation}
There are two complex odd components, or four real odd directions. A superfield is a function on this superspace:
\begin{equation}
\Phi=\Phi(x,\theta,\bar\theta).
\end{equation}

These are not tiny spatial dimensions in the metric sense. A Grassmann coordinate has no positive numerical magnitude. Calling it small is a mnemonic for nilpotence:
\begin{equation}
(\theta^\alpha)^2=0,
\qquad
\theta^\alpha\theta^\beta=-\theta^\beta\theta^\alpha.
\end{equation}
\editorialnote{Nilpotence, not numerical smallness, is the reason every expansion in finitely many Grassmann coordinates terminates.}

With four independent odd coordinates, a completely general superfield has at most \(2^4=16\) Grassmann monomials. Schematically,
\begin{equation}
\begin{aligned}
\Phi(x,\Theta)={}&\phi(x)+\Theta_A\psi_A(x)\\
&+\frac{1}{2}\Theta_A\Theta_B Z_{AB}(x)+\cdots\\
&+\Theta_1\Theta_2\Theta_3\Theta_4D(x),
\end{aligned}
\end{equation}
where \(\Theta_A\) collectively denotes the four real odd coordinates. Since the whole superfield may be chosen even, coefficients of even powers are bosonic and coefficients of odd powers are fermionic. One superfield therefore packages a finite multiplet of ordinary component fields.

This packaging is more than notation. Supersymmetry acts geometrically as translation in the odd directions mixed with spacetime translation. A Lagrangian constructed from superfields according to superspace rules inherits the symmetry in the same way that a scalar Lagrangian constructed from tensors inherits rotational or Lorentz invariance.

The action now includes Grassmann measures. A typical full-superspace term has the form
\begin{equation}
S_D=\int d^4x\,d^2\theta\,d^2\bar\theta\,
\mathcal K(\Phi,\Phi^\dagger).
\end{equation}
Chiral terms instead use
\begin{equation}
S_F=\int d^4x\,d^2\theta\,W(\Phi)+\mathrm{h.c.}
\end{equation}
\editorialnote{The distinction between full-superspace and chiral measures is the standard refinement of the schematic \(d^4x\,d^4\theta\) action. Concrete examples will determine which measure is appropriate.}

We have therefore reached a question that ordinary calculus does not answer: what does integration over an anticommuting coordinate mean?

\section{Berezin Integration}

\subsection{The ordinary model: a total derivative}

Begin with an ordinary function \(F(x)\) that vanishes at large positive and negative \(x\). The fundamental theorem of calculus gives
\begin{equation}
\int_{-\infty}^{\infty}dx\,\frac{dF}{dx}
=F(\infty)-F(-\infty)
=0.
\end{equation}

\begin{figure}[H]
\centering
\includegraphics[width=0.88\textwidth]{lecture_06_figure_05.png}
\lectureframe{01:26:31.594}
\caption{The ordinary prototype for Berezin integration: an integral of a total derivative reduces to endpoint values and vanishes when the function dies off. The lower board line is partly obscured, so the clean equation is supplied in the text.}
\end{figure}

We want the Grassmann integral to be linear and to share the useful property that a total derivative integrates to zero. For one odd coordinate, require
\begin{align}
\int d\theta\,[af(\theta)+bg(\theta)]
&=
a\int d\theta\,f(\theta)
+b\int d\theta\,g(\theta),\\
\int d\theta\,\partial_\theta f(\theta)&=0.
\end{align}
The second condition is the algebraic analogue of translation invariance of the measure. There are no endpoints of an odd number line and no sum over tiny intervals.

\begin{classroomqa}
\audiencequestion{If there is no indefinite Grassmann integral, are we abandoning the idea that integration is the inverse of differentiation? \lecturetimestamp{01:21:19}}
\lecturerresponse{Yes. Berezin integration is a definite linear functional, not an antiderivative operation. For one Grassmann variable it is, in fact, represented by the same algebraic operation as differentiation.}
\end{classroomqa}

\begin{classroomqa}
\audiencequestion{Is the \(d^4x\) integration over all spacetime also a definite integral? \lecturetimestamp{01:21:34}}
\lecturerresponse{Yes. Boundary conditions are understood when one integrates the action over spacetime. The difference is that ordinary spacetime integration can be understood through limits of sums, whereas Grassmann integration is defined algebraically by its action on the finite basis \(1,\theta\).}
\end{classroomqa}

\subsection{One variable}

Every function of one odd variable can be written with its coefficient in the fixed right-hand order
\begin{equation}
f(\theta)=a+\theta b.
\end{equation}
Since \(\partial_\theta\theta=1\), the total-derivative rule implies
\begin{equation}
\int d\theta\,1=0.
\end{equation}
The remaining normalization is chosen to be
\begin{equation}
\int d\theta\,\theta=1.
\end{equation}
Consequently,
\begin{equation}
\int d\theta\,f(\theta)=b.
\end{equation}
The integral selects the coefficient of the highest available power. With the left-derivative convention used here,
\begin{equation}
\partial_\theta f=b
=\int d\theta\,f.
\end{equation}

\begin{classroomqa}
\audiencequestion{Is there a function \(F(\theta)\) satisfying \(\partial_\theta F=\theta\)? \lecturetimestamp{01:30:05}}
\lecturerresponse{No. Such a function would need a term proportional to \(\theta^2\), but \(\theta^2=0\). This is another way to see why there is no Grassmann analogue of an indefinite antiderivative.}
\end{classroomqa}

\subsection{Two variables and the order of the measure}

For a complex pair \(\theta,\bar\theta\), write a general even function in the ordered form
\begin{equation}
f(\theta,\bar\theta)
=a+\bar\theta\psi+\bar\psi\theta+b\,\bar\theta\theta.
\end{equation}
No higher term exists because
\begin{equation}
\theta^2=\bar\theta^2=0,
\qquad
\bar\theta\theta=-\theta\bar\theta.
\end{equation}

We now fix the convention
\begin{equation}
\int d\theta\,d\bar\theta\,
\bar\theta\theta=1.
\end{equation}
With iterated integrals, the adjacent \(d\bar\theta\) operation acts first. The table is then
\begin{align}
\int d\theta\,d\bar\theta\,1&=0,\\
\int d\theta\,d\bar\theta\,\theta&=0,\\
\int d\theta\,d\bar\theta\,\bar\theta&=0,\\
\int d\theta\,d\bar\theta\,\bar\theta\theta&=1.
\end{align}
Thus
\begin{equation}
\int d\theta\,d\bar\theta\,
f(\theta,\bar\theta)=b.
\end{equation}

\begin{classroomqa}
\audiencequestion{For a term linear in \(\theta\), is it really the \(\theta\) integration that makes the double integral vanish? \lecturetimestamp{01:34:59}}
\lecturerresponse{No; the live explanation had the roles reversed. Integrating \(\theta\) with respect to \(\theta\) gives one. The result then contains no \(\bar\theta\), and the remaining \(\bar\theta\) integral of one vanishes. Either way, a double integral survives only when both variables are present.}
\end{classroomqa}

\begin{classroomqa}
\audiencequestion{If \(\theta\) and \(\bar\theta\) are interchanged, does the integral change sign? \lecturetimestamp{01:36:33}}
\lecturerresponse{Yes. Odd variables and odd measures anticommute. With the convention above,
\[
\int d\bar\theta\,d\theta\,\bar\theta\theta=-1.
\]
One must therefore choose an ordering and keep it consistently. The physics is convention-independent, but intermediate signs are not.}
\end{classroomqa}

For \(n\) independent Grassmann variables, the pattern is immediate:
\begin{equation}
\int d\theta_n\cdots d\theta_1\,
f(\theta_1,\ldots,\theta_n)
\end{equation}
extracts the coefficient of the ordered top monomial
\(\theta_1\cdots\theta_n\), up to the sign fixed by the chosen order of variables and measures. This simple component-extraction rule is the engine behind superspace actions. A supersymmetry variation becomes an odd total derivative, and the superspace integral discards it, just as an ordinary action discards a spacetime total derivative under suitable boundary conditions.

\section{Why Such a Small Calculus Has So Much Power}

The rules look almost trivial:
\begin{equation}
\int d\theta\,1=0,
\qquad
\int d\theta\,\theta=1.
\end{equation}
Yet they let us build interacting field theories in which relations between bosonic and fermionic masses and couplings are preserved automatically. The purpose of superspace calculus is the same as the purpose of tensor calculus in relativity: arrange the notation so that the desired symmetry survives every allowed manipulation.

Sign bookkeeping is the price. Every exchange of two odd objects contributes a minus sign, and a calculation can be conceptually correct while failing by one ordering choice. That difficulty is not decorative; it is the algebraic shadow of fermionic statistics.

\begin{classroomqa}
\audiencequestion{Were Grassmann numbers invented together with supersymmetry, or did they exist earlier? \lecturetimestamp{01:41:19}}
\lecturerresponse{They are much older. Hermann Grassmann developed the underlying exterior algebra in the nineteenth century. Differential forms already use anticommuting basis elements; twentieth-century quantum field theory and supersymmetry found a new physical role for the same mathematics.}
\end{classroomqa}

We are now ready for concrete constructions. Superfields collect complete boson-fermion multiplets, superspace derivatives preserve the algebra, and Berezin integration extracts the component action. The next task is to use this machinery in the simplest one-dimensional model and then carry it into field theory.
